\subsection{SEMI-CONTINUOUS FUNCTIONS ON A UNIFORMIZABLE SPACE}

In Chapter IV, § 6, no. 2, Corollary to Theorem 4, we showed that the upper envelope of a family of continuous real-valued functions on a topological space is a lower semi-continuous function. If the space is uniformizable, there is a converse to this proposition:

PROPOSITION 5. In order that every lower semi-continuous real-valued function \(f\) (finite or not) on a topological space \(\mathbf{X}\) should be the upper envelope of the continuous real-valued functions on \(\mathbf{X}\) (finite or not) which are \(\leqslant f\) , it is necessary and sufficient that \(\mathbf{X}\) be uniformizable.

The condition is necessary. Let \(x_0\) be any point of \(\mathbf{X}\) and let \(\mathbf{V}\) be any open neighbourhood of \(x_0\) ; then the characteristic function \(\varphi_{\mathbf{V}}\) of the set \(\mathbf{V}\) is lower semi-continuous (Chapter IV, § 6, no. 2, Corollary to Proposition 1); by hypothesis, there is therefore a continuous real-valued function \(g\) on \(\mathbf{X}\) such that \(g \leqslant \varphi_{\mathbf{V}}\) and \(g(x_0) = a > 0\) . The continuous function \(\inf\left(\mathbf{i}, \frac{\mathbf{i}}{a} g^+\right)\) takes its values in {[}o, i{]}, is equal to o in CV, and equal to i at \(x_0\) . Hence (Theorem 2) \(\mathbf{X}\) is uniformizable.

The condition is sufficient. Consider first the case in which \(f\) takes its values in \([-1, +1]\) . We have to show that, for each \(x_0 \in X\) and each number \(a < f(x_0)\) , there is a continuous real-valued function \(g\) on \(X\) such that \(g \leqslant f\) and \(g(x_0) \geqslant a\) . If \(a \leqslant -1\) , we may take \(g\) to be the constant \(-1\) . If \(-1 < a < f(x_0)\) , there is a neighbourhood \(V\) of \(x_0\) such that \(f(x) \geqslant a\) for all \(x \in V\) . Since \(X\) is uniformizable, there is a continuous real-valued function \(h\) on \(X\) , with values in \([0, 1]\) , such that \(h(x_0) = 0\) and \(h(x) = 1\) for all \(x \in C_V\) . We may then take \(g(x) = a - (a + 1)h(x)\) , and we have a continuous function satisfying the stated conditions. Note that this function takes its values in \([-1, +1]\) .

The general case follows by transfer of structure; for there is a strictly increasing homeomorphism of \([-1, +1]\) onto \(\overline{R}\) (Chapter IV, § 4, no. 2, Proposition 2), and the definition of a semi-continuous function involves only the order structure and the topology of \(\overline{R}\) .

Remark. In the above proof we see that the function g does not take the value +1. By transfer of structure it follows that every lower semi-continuous real-valued function f on the uniformizable space X is the upper envelope of the continuous real-valued functions \(g \leqslant f\) on X which do not take the value +∞.

